%%%%%%%%%%%%%%%%%%%%%%%%%%%%%%%%%%%%%%%%%%%%%%%%%%%%%%%
%% Preâmbulo do template UFC - Pós-Graduação         %%
%%                                                   %%
%% Adaptado a partir do template original UFC para   %%
%% trabalhos de disciplina de pós-graduação com foco %%
%% em Engenharia Elétrica / Eletrônica de Potência.  %%
%%%%%%%%%%%%%%%%%%%%%%%%%%%%%%%%%%%%%%%%%%%%%%%%%%%%%%%

% =====================================================
%   Codificação, idioma e tipografia
% =====================================================
\usepackage[utf8]{inputenc}                         % Acentuação direta
\usepackage[T1]{fontenc}                            % Codificação da fonte em 8 bits
\usepackage{microtype}                              % Melhorias finas de justificação
\usepackage{indentfirst}                            % Endenta o primeiro parágrafo de cada seção
\usepackage{mathptmx}                               % Fonte Times New Roman para texto e matemática

% =====================================================
%   Matemática
% =====================================================
\usepackage{amsfonts, amssymb, amsmath}             % Símbolos e ambientes matemáticos
\usepackage{amsgen}
\usepackage{mathtools}                              % Extensões de amsmath (matrizes, alinhamentos)

% Operadores em português (não há \sen, \tg, \cotg em LaTeX padrão)
\DeclareMathOperator{\sen}{sen}
\DeclareMathOperator{\tg}{tg}
\DeclareMathOperator{\cotg}{cotg}
\DeclareMathOperator{\arctg}{arctg}

% =====================================================
%   Figuras, tabelas e legendas
%   (memoir/abntex2 já fornecem [H] para floats e
%   ambiente subcaption via \newsubfloat — por isso
%   NÃO carregamos float nem subcaption aqui.)
% =====================================================
\usepackage{graphicx}                               % Inserir figuras
\graphicspath{{figuras/}{figuras/circuitos/}{figuras/formas-de-onda/}}
\usepackage{booktabs}                               % Tabelas profissionais (\toprule etc.)
\usepackage{multirow, array}                        % Múltiplas linhas e colunas em tabelas
\usepackage[font=singlespacing,format=plain,justification=justified,skip=0pt,singlelinecheck=false]{caption}

% =====================================================
%   Unidades SI e notação numérica (siunitx)
% =====================================================
\usepackage{siunitx}
\sisetup{
	output-decimal-marker = {,},
	per-mode               = symbol,
	inter-unit-product     = \ensuremath{{}\cdot{}},
	exponent-product       = \cdot,
	separate-uncertainty   = true,
	range-units            = single,
	range-phrase           = {\ \text{a}\ },
	list-units             = single,
	list-pair-separator    = {\ \text{e}\ },
	list-final-separator   = {\ \text{e}\ }
}

% =====================================================
%   Diagramas de circuitos (CircuiTikZ)
% =====================================================
\usepackage{tikz}
\usepackage[siunitx,european,straightvoltages,RPvoltages]{circuitikz}
\usetikzlibrary{calc,arrows.meta,positioning,patterns,decorations.markings}
\ctikzset{
	bipoles/length=1.0cm,
	diodes/scale=0.7,
	transistors/scale=0.8,
	resistors/scale=0.7,
	capacitors/scale=0.7,
	inductors/scale=0.7,
	sources/scale=0.8
}

% =====================================================
%   Gráficos vetoriais (pgfplots) — formas de onda,
%   curvas de FP e THD em função de Po, etc.
% =====================================================
\usepackage{pgfplots}
\pgfplotsset{
	compat=1.18,
	every axis/.append style={
		font=\small,
		line width=0.6pt,
		axis lines=left,
		grid=both,
		major grid style={line width=.2pt,draw=gray!30},
		minor grid style={line width=.1pt,draw=gray!15},
		tick align=outside
	},
	every axis plot/.append style={line width=1pt}
}
\usepgfplotslibrary{groupplots,fillbetween}

% =====================================================
%   Listagens de código e algoritmos
% =====================================================
\usepackage{verbatim}
\usepackage[svgnames,table]{xcolor}                 % Cores extras (Green, Blue, etc.)
\usepackage{listings}
\usepackage[portuguese,ruled,lined]{algorithm2e}
\usepackage{algorithmic}

% =====================================================
%   Estrutura, formatação e auxiliares
% =====================================================
\usepackage{tocloft}                                % Personaliza o Sumário
\usepackage{etoolbox}                               % Toolbox de comandos
\usepackage[nogroupskip,nonumberlist]{glossaries}   % Lista de abreviaturas e siglas
\usepackage[alf, abnt-emphasize=bf, recuo=0cm, abnt-etal-cite=2, abnt-etal-list=0, abnt-etal-text=it]{abntex2cite}
\usepackage{lipsum}                                 % Texto fictício (Lorem Ipsum)
\usepackage{appendix}                               % Geração de apêndices
\usepackage{paracol}                                % Parágrafos sem identação
\usepackage{pdfpages}                               % Inclui PDFs externos
\usepackage{lib/ufctex}                             % Normas UFC

% =====================================================
%   Hyperref — já é carregado pelo abntex2; apenas
%   complementamos a configuração para melhor SyncTeX
%   e bookmarks bem ordenados no leitor de PDF.
% =====================================================
\hypersetup{
	bookmarksopen        = true,
	bookmarksnumbered    = true,
	bookmarksopenlevel   = 1,
	unicode              = true,
	pdfencoding          = auto,
	pdfstartview         = FitH,
	breaklinks           = true,
	colorlinks           = false
}

% =====================================================
%   Glossário e índice
% =====================================================
\makeglossaries
\makeindex
